On introduit $\qty{3.27}{\g}$ de poudre de zinc dans $\qty{100}{\mL}$ d'une
solution de sulfate de cuivre(II) de concentration molaire
$c=\qty{1.0}{\mol\per\L}$. La réaction produit du cuivre métal et des ions
\ch{Zn^{2+}_{(aq)}}. Les ions sulfate n'interviennent pas.
\begin{enumerate}
    \item Écrire l'équation de la dissolution dans l'eau du sulfate de
          cuivre(II). En déduire les concentrations
          $\bigl[\ch{Cu^{2+}_{(aq)}}\bigr]$ et
          $\bigl[\ch{SO_{4}^{2-}_{(aq)}}\bigr]$ dans la solution.
    \item Écrire l'équation de la réaction entre le zinc et les ions
          \ch{Cu^{2+}_{(aq)}}.
    \item Calculer les quantités de matière initiales des réactifs.
    \item Quel est le réactif limitant~? Que vaut l'avancement maximal~?
    \item En déduire, dans l'état final~:
          \begin{itemize}
            \item la composition en quantité de matière du système~;
            \item la masse de cuivre obtenue~;
            \item la concentration molaire des ions \ch{Zn^{2+}_{(aq)}}~;
            \item la concentration molaire des ions \ch{Cu^{2+}_{(aq)}}.
          \end{itemize}

          On supposera que, pendant la réaction, le volume de la solution ne
          varie pas.
\end{enumerate}

